%% titre: Devoir surveillé — Optimisation
%% niveau: Licence 3
%% matiere: Optimisation
%% annee: 2025-2026
%% description: Devoir d'Optimisation licence 3 mathématique. Durée : 2 h, sans documents.

	\begin{center}
		\textbf{\sc Examen d'Optimisation}\\
		Filière : L3S5 Mathématiques\\
		Durée : 2h00mn\\
	\end{center}
\qquad\\
\textbf{Exercice 1} (04 points):\\

\begin{enumerate}
	\item Soit $f$ une fonction différentiable sur un ensemble convexe 
	$C \subset \mathbb{R}^n$.  
	En utilisant la définition de la convexité et de la dérivée directionnelle, 
	montrer que si $f$ est convexe alors :
	\[
	f(y) \ge f(x) + \langle \nabla f(x), y - x \rangle, \quad \forall x,y \in C.
	\]
	
	\item Déduire que $\nabla f$ est monotone sur $C$, c’est-à-dire :
	\[
	\langle \nabla f(y) - \nabla f(x),\, y - x \rangle \ge 0, \quad \forall x,y \in C.
	\]
\end{enumerate}

\qquad\\
\textbf{Exercice 2} (06 points):\\

Soit $\alpha \in \mathbb{R}$.  
On considère la fonction $f_\alpha : \mathbb{R}^2 \to \mathbb{R}$ définie par :
\[
f_\alpha(x,y) = x^2 + y^2 + \alpha xy + x + 2y.
\]

\begin{enumerate}
	\item Pour quelles valeurs de $\alpha$ la fonction $f_\alpha$ est-elle convexe ?  
	Et strictement convexe ?
	
	\item Calculer les points critiques de $f_\alpha$ en fonction de $\alpha$ puis
	résoudre le problème :
	\[
	\min_{(x,y)\in\mathbb{R}^2} f_\alpha (x,y)
	\quad \text{lorsque } \alpha \ne \pm 2.
	\]
\end{enumerate}

\qquad\\
\textbf{Exercice 3} (04 points):\\

Une entreprise veut déménager son matériel composé de 450 machines 
de trois types : $M_1$, $M_2$, $M_3$.  
Elle décide de louer des camions.  
La société de location dispose de trois sortes de véhicules : 
$V_1$, $V_2$, $V_3$, dont les tarifs sont respectivement de  
50000, 80000 et 120000 FCFA par voyage.

Les camions $V_1$ peuvent chacun transporter :  
1 machine $M_1$, 4 machines $M_2$, et 10 machines $M_3$.

Pour des raisons techniques la place d’une machine d’un type donné ne peut pas être utilisée pour une machine d’un autre type.

Chaque camion $V_2$ peut transporter :  
2 machines $M_1$, 6 machines $M_2$, 20 machines $M_3$.

Alors que les véhicules $V_3$ ont pour capacité maximale :  
4 machines $M_1$, 20 machines $M_2$, 24 machines $M_3$.

On veut transporter en un seul convoi :  
30 machines $M_1$, 120 machines $M_2$, et 300 machines $M_3$.

L’entreprise veut déterminer le nombre de véhicules à louer pour minimiser le coût total de transport.  
Donner le modèle linéaire de ce problème sans le résoudre.

\qquad\\
\textbf{Exercice 4} (06 points):\\

Résoudre le programme linéaire suivant par la méthode du simplexe :

\[
\max\; Z = 3x_1 + 5x_2 + 4x_3
\]

Sous contraintes :
\[
\begin{cases}
2x_1 + x_2 + x_3 \le 10, \\
x_1 + 3x_2 + 2x_3 \le 11, \\
x_1 + x_2 + 4x_3 \le 12, \\
x_1, x_2, x_3 \ge 0.
\end{cases}
\]
